\documentclass[10pt,a4paper]{amsart}
\usepackage[margin=19mm]{geometry}
\usepackage{amsmath,amssymb,mathtools}
\usepackage[scr=rsfs,cal=euler]{mathalfa}
\usepackage{xfrac}
\usepackage[unicode,hidelinks,bookmarks=false]{hyperref}
\newcommand{\ie}{i.e.\ }
\newcommand{\cf}{cf.\ }
\newcommand{\resp}{respectively\ }
\newcommand{\Subsection}{\subsection}
\providecommand{\nobreakdash}{\nobreak\hbox{-}}
\setlength{\emergencystretch}{2em}
\multlinegap0pt
\usepackage{bm}
\usepackage{bbm}
\usepackage{dsfont}
\usepackage{xspace}
\usepackage{cancel}
\usepackage{mathtools}
\usepackage{amsthm}
\makeatletter
\@ifundefined{thm}{\newtheorem{thm}{Theorem}}{}
\makeatother
\def\bH{\mbox{\boldmath $H$}}
\def\bK{\mbox{\boldmath $K$}}
\def\br{\mbox{\boldmath $r$}}
\def\uM{\underline{M}}
\title[On the governing equations of membrane O surfaces]{On the governing equations of membrane O surfaces}
\author{Yoshiki Jikumaru}
\date{Source: arXiv:2510.19284v1 (CC BY 4.0)}
\begin{document}
\maketitle

\noindent\emph{Source and licence.} On the governing equations of membrane O surfaces. Version arXiv:2510.19284v1 (CC BY 4.0), distributed under the Creative Commons Attribution 4.0 International licence, as recorded in the arXiv licence metadata for this exact version and shown on the abstract page https://arxiv.org/abs/2510.19284v1. Full rights notice: licenses/arxiv-2510.19284-CC-BY-4.0.txt.\par\smallskip

\noindent\textit{Editorial scope.} Counted unit: the Thm whose source label is \texttt{thm:main2\_2nd} in the article (source0.tex lines 765--784, source label \texttt{thm:main2\_2nd}), delivered together with the article's own complete original proof of it (source0.tex lines 786--859, one \texttt{proof} environment of 74 lines). No mathematics is written, repaired or re-derived: statement and proof are the article's own text line for line. Required context retained uncounted: eq:1st\_govern (lines 335--345); eq:2nd\_govern (lines 372--380); eq:Lax\_pair (lines 662--709); eq:eqn\_H\_K (lines 723--731); thm:main2\_1st (lines 744--763). Every definition, display and label of those lines is retained unchanged. Omitted from this package and not redistributed: 1--334, 346--371, 381--661, 710--722, 732--743, 764--764, 785--785, 860--973. Nothing inside a retained excerpt is omitted or altered: every retained body is delivered exactly as the article's lines are written, so no omission receipt of any kind accompanies this package. Citations: the retained excerpts carry no citation command at all, so no bibliography entry of the article is transcribed and no reference is redistributed. Reference adaptation: none was needed and none was made; every reference inside the delivered unit resolves inside it, so no unresolved reference or citation is produced. Printed slips kept literal: no printed word, symbol, hypothesis or number of the counted statement or of its proof was altered. Layout and class: the article's own class is \texttt{amsart} with packages including amsmath, amssymb, amsthm, graphicx, color, comment, mathtools; the wrapper is the lane's pinned \texttt{amsart} editorial wrapper at 10pt on A4, which restates the theorem-like environments the retained text needs and adds the article's own macro definitions that the retained text needs (\texttt{\textbackslash bH}, \texttt{\textbackslash bK}, \texttt{\textbackslash br}, \texttt{\textbackslash uM}), with extra wrapper packages (bm, bbm, dsfont, xspace, cancel, mathtools). The delivered numbering is the unit's own. Assets: no retained span contains a figure environment, a graphics-inclusion command, a PDF-page inclusion, a table or a TikZ picture; no image file is copied into this package, the package declares no asset dependency and no image file is redistributed with it. Licence: version arXiv:2510.19284v1 of this article is distributed under the Creative Commons Attribution 4.0 International licence, as recorded in the arXiv licence metadata for this exact version and shown on the abstract page https://arxiv.org/abs/2510.19284v1. A full rights notice is delivered at licenses/arxiv-2510.19284-CC-BY-4.0.txt. Characters: every retained excerpt is pure ASCII, so the delivered engine, XeLaTeX with its default Latin Modern text font, reports no missing character.
\begin{equation}
\label{eq:1st_govern}
\begin{aligned}
&h_x = (h + \coth \alpha) \xi_x, \quad 
h_y = (h + \tanh \alpha) \xi_y,\\
&\xi_{xy} = \xi_x \xi_y + (\log \sinh \alpha)_y \xi_x + (\log \cosh \alpha)_x \xi_y,\\
&(\alpha_x + \xi_x \coth \alpha)_x 
+ (\alpha_y + \xi_y \tanh \alpha)_y
+ e^{2\xi} \sinh \alpha \cosh \alpha = 0.
\end{aligned}
\end{equation}

\begin{equation}
\label{eq:2nd_govern}
\begin{aligned}
&h_x = (h + \cot \alpha) \xi_x, \quad
h_y = (h - \tan \alpha) \xi_y,\\
&\xi_{xy} = \xi_x \xi_y + (\log \sin \alpha)_y \xi_x + (\log \cos \alpha)_x \xi_y,\\
&(-\alpha_x + \xi_x \cot \alpha)_x + (\alpha_y + \xi_y \tan \alpha)_y + e^{2\xi} \sin \alpha \cos \alpha = 0.
\end{aligned}
\end{equation}

\begin{equation}
\label{eq:Lax_pair}
\begin{aligned}
\begin{pmatrix}
\lambda \\
\mu \\
\omega \\
\varphi \\
\chi
\end{pmatrix}_x &=
\begin{pmatrix}
0 & -p & m \overline{A}_1 - H_\circ & - m q_n A_1 & m H_\circ \\
p & 0 & 0 & 0 & 0 \\
H_\circ & 0 & 0 & 0 & 0 \\
A_1 & 0 & 0 & 0 & 0 \\
\overline{A}_1 & 0 & 0 & 0 & 0
\end{pmatrix}
\begin{pmatrix}
\lambda \\
\mu \\
\omega \\
\varphi \\
\chi
\end{pmatrix}, \\
%
\begin{pmatrix}
\lambda \\
\mu \\
\omega \\
\varphi \\
\chi
\end{pmatrix}_y &=
\begin{pmatrix}
0 & q & 0 & 0 & 0 \\
-q & 0 & m \overline{A}_2 - K_\circ & - m q_n A_2 & m K_\circ \\
0 & K_\circ & 0 & 0 & 0 \\
0 & A_2 & 0 & 0 & 0 \\
0 & \overline{A}_2 & 0 & 0 & 0
\end{pmatrix}
\begin{pmatrix}
\lambda \\
\mu \\
\omega \\
\varphi \\
\chi
\end{pmatrix}.
\end{aligned}
\end{equation}

\begin{equation}
\label{eq:eqn_H_K}
\begin{aligned}
H &= m \bH \Lambda \uM^T
= m \left( \nu H_\circ + \omega \overline{A}_1 - q_n \varphi A_1 + \frac{q_n \varphi^2}{2\omega} H_\circ \right),\\
K &= m \bK \Lambda \uM^T
= m \left( \nu K_\circ + \omega \overline{A}_2 - q_n \varphi A_2 + \frac{q_n \varphi^2}{2\omega} K_\circ \right).
\end{aligned}
\end{equation}

\begin{thm}
\label{thm:main2_1st}
For a membrane O surface of the 1st kind, by using solutions of the linear system \eqref{eq:Lax_pair}, we introduce variables $\xi'$, $\alpha'$ and $h'$ as follows:
\begin{equation}
e^{\xi'} = \frac{q_n}{2} \frac{\omega}{\nu} e^{-\xi} (1-t^2), \quad
e^{\alpha'} = e^{-\alpha} \frac{1-t}{1+t}, \quad
h' = t + \frac{\varphi}{\omega} e^{\xi'}.
\end{equation}
Here we put $t = h - \dfrac{\varphi}{\omega} e^\xi$.
Then, the coefficients of the first and third fundamental forms of $\br'$ are given by
\begin{equation}
\begin{aligned}
A_1' &= \cosh \alpha' + h' \sinh \alpha', \quad
A_2' = - (\sinh \alpha' + h' \cosh \alpha'),\\
H_\circ' &= e^{\xi'} \sinh \alpha', \quad
K_\circ' = - e^{\xi'} \cosh \alpha'.
\end{aligned}
\end{equation}
Therefore, $(\xi', \alpha', h')$ again satisfies the equation \eqref{eq:1st_govern}, and the B\"acklund transformation for membrane O surfaces preserves the membrane O surface of the 1st kind.
\end{thm}

\begin{thm}
\label{thm:main2_2nd}
For a membrane O surface of the 2nd kind, by using solutions of the linear system \eqref{eq:Lax_pair}, we introduce variables $\xi'$, $\alpha'$ and $h'$ as follows:
\begin{equation}
e^{\xi'} = \frac{q_n}{2} \frac{\omega}{\nu} e^{-\xi} (1+t^2), \quad
e^{i\alpha'} = e^{i\alpha} \frac{1-it}{1+it}, \quad
h' = -t + \frac{\varphi}{\omega}e^{\xi'}.
\end{equation}
Here we put $t = h - \dfrac{\varphi}{\omega} e^\xi$.
Then, the coefficients of the first and third fundamental forms of $\br'$ are given by
\begin{equation}
\begin{aligned}
A_1' &= \cos \alpha' + h' \sin \alpha', \quad
A_2' = \sin \alpha' - h' \cosh \alpha',\\
H_\circ' &= e^{\xi'} \sin \alpha', \quad
K_\circ' = - e^{\xi'} \cos \alpha'.
\end{aligned}
\end{equation}
Therefore, $(\xi', \alpha', h')$ again satisfies equation \eqref{eq:2nd_govern}, and the B\"acklund transformation for membrane O surfaces preserves the membrane O surface of the 2nd kind.
\end{thm}

\begin{proof}[Proof of Theorem~\ref{thm:main2_1st}]
The coefficients for the fundamental form in the 1st kind are restated as follows:
\begin{equation}
\begin{aligned}
A_1 &= \cosh \alpha + h \sinh \alpha, \quad
A_2 = \sinh \alpha + h \cosh \alpha,\\
H_\circ &= e^\xi \sinh \alpha, \quad
K_\circ = e^\xi \cosh \alpha,\\
\overline{A}_1 &= \frac{q_n}{2} e^{-\xi} ((1+h^2)\sinh \alpha + 2h \cosh \alpha),\\
\overline{A}_2 &= \frac{q_n}{2} e^{-\xi} ((1+h^2)\cosh \alpha + 2h \sinh \alpha).
\end{aligned}
\end{equation}
Under these conditions, setting $t = h - \dfrac{\varphi}{\omega} e^\xi$ yields the following equation:
\begin{equation}
\begin{aligned}
\omega \overline{A}_1 - q_n \varphi A_1 + \frac{q_n \varphi^2}{2\omega} H_\circ
&= \frac{q_n}{2} \omega e^{-\xi} ((1+t^2) \sinh \alpha + 2t \cosh \alpha),\\
\omega \overline{A}_2 - q_n \varphi A_2 + \frac{q_n \varphi^2}{2\omega} K_\circ
&= \frac{q_n}{2} \omega e^{-\xi} ((1+t^2) \cosh \alpha + 2t \sinh \alpha).
\end{aligned}
\end{equation}
Therefore, if we put
\begin{equation}
e^{\xi'} = \frac{q_n}{2} \frac{\omega}{\nu} e^{-\xi} (1-t^2), \quad
e^{\alpha'} = e^{-\alpha} \frac{1-t}{1+t},
\end{equation}
then we have
\begin{equation}
\begin{aligned}
\omega \overline{A}_1 - q_n \varphi A_1 + \frac{q_n \varphi^2}{2\omega} H_\circ
&= -\nu e^{\xi'} \sinh \alpha',\\
\omega \overline{A}_2 - q_n \varphi A_2 + \frac{q_n \varphi^2}{2\omega} K_\circ
&= \nu e^{\xi'} \cosh \alpha'.
\end{aligned}
\end{equation}
Using these, equation \eqref{eq:eqn_H_K} can be rewritten as follows:
\begin{equation}
H = m\nu (H_\circ - e^{\xi'} \sinh \alpha'), \quad
K = m\nu (K_\circ + e^{\xi'} \sinh \alpha'),
\end{equation}
Therefore, the coefficients $H_\circ'$ and $K_\circ'$ of the third fundamental form of $\br'$ are given by
\begin{equation}
H_\circ' 
= H_\circ - \frac{H}{m\nu} = e^{\xi'} \sinh \alpha',\quad
K_\circ' 
= K_\circ - \frac{K}{m\nu} = - e^{\xi'} \cosh \alpha'.
\end{equation}
Furthermore, $h'$ is defined by the following relation:
\begin{equation}
h' 
= t + \frac{\varphi}{\omega}e^{\xi'}
= h + \frac{\varphi}{\omega}(e^{\xi'} - e^\xi).
\end{equation}
Then, the coefficients $A_1'$ and $A_2'$ of the first fundamental form of $\br'$ are given by the following expressions:
\begin{equation}
\begin{aligned}
A_1' &= A_1 - \frac{H}{m \omega \nu} \varphi
= \cosh \alpha + t \sinh \alpha + \frac{\varphi}{\omega} e^{\xi'} \sinh \alpha'
= \cosh \alpha' + h' \sinh \alpha',\\
A_2' &= A_2 - \frac{K}{m \omega \nu} \varphi
= \sinh \alpha + t \sinh \alpha - \frac{\varphi}{\omega} e^{\xi'} \cosh \alpha'
= - (\sinh \alpha' + h' \cosh \alpha').
\end{aligned}
\end{equation}
Here, we used the relations
\begin{equation}
\begin{aligned}
\cosh \alpha + t \sinh \alpha &= \cosh \alpha_1 + t \sinh \alpha_1, \\
\sinh \alpha + t \cosh \alpha &= - (\sinh \alpha_1 + t \cosh \alpha_1).
\end{aligned}
\end{equation}
This proves the theorem.
The proof for Theorem~\ref{thm:main2_2nd} is exactly the same.
\end{proof}

\end{document}
